\chapter{Threads: the internal pool}
\label{ch:parallel}

By default BigMath runs a small pool of worker threads and uses it \emph{inside} single
operations: one large multiplication is spread over several cores. You do not have to do
anything to use it, and you can switch it off at build time
(\code{-DBIGMATH\_USE\_THREADS=0}). The code is in \file{src/common/Parallel.cpp}.

\section{What runs in parallel}

\begin{center}
\small
\begin{tabularx}{\textwidth}{lX}
\toprule
operation & how the work is split \\
\midrule
addition / subtraction & if enabled by the profile, equal-length magnitudes use blocks
                         with a carry or borrow fixup (\cref{sec:paralleladdsub}) \\
NTT product & Ordinary CRT batches operand--prime transforms. The fused matrix method
              schedules transpose tiles across the pool at each stage. Finalization
              splits the carries into ranges (\cref{sec:nttthreads}). \\
parse / print $\ge$ 100\,000 digits & 8 independent subtrees of the divide-and-conquer tree \\
division, conversion & indirectly: their large products are NTT products \\
Miller--Rabin & independent randomized witnesses once base 2 passes \\
next probable prime & base-2 screens of sieve survivors, then randomized tests in
                      candidate order (\cref{sec:prime}) \\
\bottomrule
\end{tabularx}
\end{center}

Small operations never touch the pool. Schoolbook, Karatsuba, Toom-3 and long division are
single-threaded: at their sizes, handing work to another core costs more than doing it.

\section{Size of the pool}

The pool has $\min(\text{hardware threads}, 8)$ threads (\code{BIGMATH\_MAX\_THREADS}),
counting the calling thread, which always does a share of the work itself. Eight is a
deliberate cap: on the M1 Max more threads competed for the shared L2 cache during NTTs and
did not help.

\section{Warm and cold}
\label{sec:warmcold}

A worker that has nothing to do spins for 100\,\textmu s (\code{BIGMATH\_POOL\_SPIN\_US}) waiting for
new work, then goes to sleep. Waking a sleeping thread takes the operating system 20--30\,\textmu s.
For a large product that is nothing. For a mid-sized NTT product, which takes about 25\,\textmu s
with all cores working, it is as long as the whole job.

\begin{figure}[ht]
\centering
\begin{tikzpicture}[x=0.11cm,y=0.5cm]
  \node[lbl,anchor=east] at (-1,0) {caller};
  \node[lbl,anchor=east] at (-1,-1) {worker};
  \node[lbl,anchor=east] at (-1,-2) {worker};
  \node[font=\small\bfseries,anchor=west] at (0,1.2) {warm pool};
  \fill[sOrange!50] (0,-0.3) rectangle (24,0.3);
  \fill[sOrange!50] (0.5,-1.3) rectangle (24,-0.7);
  \fill[sOrange!50] (0.5,-2.3) rectangle (24,-1.7);
  \draw[ink2] (24.5,0.5) -- (24.5,-2.5) node[below,lbl] {done $\approx$ 24\,\textmu s};
  \begin{scope}[shift={(48,0)}]
    \node[font=\small\bfseries,anchor=west] at (0,1.2) {cold pool};
    \fill[sOrange!50] (0,-0.3) rectangle (50,0.3);
    \fill[rule] (0,-1.3) rectangle (25,-0.7); \node[font=\tiny] at (12.5,-1) {waking up};
    \fill[sOrange!30] (25,-1.3) rectangle (50,-0.7);
    \fill[rule] (0,-2.3) rectangle (28,-1.7);
    \fill[sOrange!30] (28,-2.3) rectangle (50,-1.7);
    \draw[ink2] (50.5,0.5) -- (50.5,-2.5) node[below,lbl] {done $\approx$ 52\,\textmu s};
  \end{scope}
\end{tikzpicture}
\caption{A small NTT product with a warm and with a cold pool (schematic). Cold, the caller
does most of the work alone while the workers wake up, so the product costs about its
single-threaded time.}
\end{figure}

This is why the multiplication dispatcher checks the pool's temperature for mid-sized
products (\cref{sec:poolrule}): when cold it uses Karatsuba, which does not need the pool,
and if cold operations arrive in a quick burst it nudges the pool awake so the rest of the
burst can use the NTT.

For parallel addition and subtraction (\cref{sec:paralleladdsub}), the warm
cutoff checks for workers that are already awake. The parked cutoff can also
start the pool. Parked benchmarks measure waking an existing pool. First use
also pays the cost of creating the workers.

\section{Nested and concurrent use}

\begin{itemize}
  \item \textbf{Nested:} a subtree worker in parallel printing may itself reach an NTT product,
        which would like to use the pool. Any parallel section started from inside a pool task
        simply runs inline on that thread. This guard is what made the 8-way conversion
        fan-out possible; an earlier attempt without it deadlocked.
  \item \textbf{Concurrent:} the pool has one job slot. If two of your threads both start large
        multiplications, the first takes the pool and the second runs its work serially on its
        own thread rather than waiting. Both are correct; the second is just not accelerated.
\end{itemize}

\section{Thread safety for your code}

\begin{itemize}
  \item Different \code{BigInteger} objects can be used from different threads freely.
  \item One object can be read from many threads at once (operators never modify their
        operands). Writing to an object while another thread uses it needs your own lock, as
        with any C++ value.
  \item A \code{CachedDivision} or \code{PreparedMultiplication} can be shared between threads
        after it is constructed; its methods are \code{const}.
  \item Internal caches (powers of ten, NTT twiddle tables, divider chains, scratch buffers)
        are \code{thread\_local}. No locking is needed, but each thread fills its own caches the
        first time. For latency-sensitive services, do one representative operation on each
        thread at startup.
\end{itemize}

\begin{keyidea}
Threads change how long an operation takes, never its result. Every algorithm choice that
depends on the pool is between two exact methods.
\end{keyidea}
